\documentclass[11pt]{amsart}
\usepackage[localtheorems]{raeez-math-template}

\providecommand{\C}{\mathbb{C}}
\providecommand{\dbar}{\bar{\partial}}
\providecommand{\g}{\mathfrak{g}}
\providecommand{\hCS}{\mathrm{hCS}}
\providecommand{\BV}{\mathrm{BV}}
\providecommand{\Sym}{\operatorname{Sym}}
\providecommand{\Obs}{\operatorname{Obs}}
\providecommand{\CoHA}{\operatorname{CoHA}}
\providecommand{\kcat}{\kappa_{\mathrm{cat}}}
\providecommand{\kch}{\kappa_{\mathrm{ch}}}
\providecommand{\kBKM}{\kappa_{\mathrm{BKM}}}
\providecommand{\kfib}{\kappa_{\mathrm{fiber}}}

\theoremstyle{plain}
\newtheorem{theorem}{Theorem}[section]
\newtheorem{proposition}[theorem]{Proposition}
\newtheorem{lemma}[theorem]{Lemma}
\theoremstyle{definition}
\newtheorem{definition}[theorem]{Definition}
\newtheorem{construction}[theorem]{Construction}
\theoremstyle{remark}
\newtheorem{remark}[theorem]{Remark}

\title{Six-dimensional holomorphic Chern--Simons on \texorpdfstring{$\C^3$}{C3}: BV data and quantisation boundaries}
\author{}
\date{}

\begin{document}
\maketitle

Let \(X\) be a complex threefold with a nowhere-vanishing holomorphic
volume form \(\Omega_X\). Let \(\g\) be a finite-dimensional Lie algebra
with invariant non-degenerate symmetric pairing
\(\langle-,-\rangle_\g\). On an open or noncompact \(X\) all pairings
below use compactly supported test fields on the left; on compact \(X\)
the support decoration may be suppressed.

\section{The local classical BV theory}

\begin{definition}[BV field space]
The holomorphic Chern--Simons BV fields are
\[
  \mathcal E_{\hCS,c}(X,\g)
  =
  \Omega^{0,\bullet}_c(X,\g)[1].
\]
The bracket is the product of the Dolbeault wedge product and the Lie
bracket of \(\g\). The superfield
\[
  \mathcal A = c + A + A^\vee + c^\vee
\]
has components
\[
\begin{array}{c|cccc}
 & c & A & A^\vee & c^\vee \\ \hline
\text{Dolbeault bidegree} & (0,0) & (0,1) & (0,2) & (0,3)\\
\text{ghost number} & 1 & 0 & -1 & -2
\end{array}
\]
and total superfield degree \(1\). The symbols \(A^\vee,c^\vee\)
denote antifields, paired respectively with \(A,c\).
\end{definition}

\begin{definition}[Classical local action]
For \(\mathcal A\in\mathcal E_{\hCS,c}(X,\g)\) set
\[
  I_{\hCS}(\mathcal A)
  =
  \int_X \Omega_X\wedge
  \left(
    \frac{1}{2}\langle \mathcal A,\dbar\mathcal A\rangle_\g
    +
    \frac{1}{6}\langle \mathcal A,[\mathcal A,\mathcal A]\rangle_\g
  \right).
\]
The displayed expression means: multiply forms, take the \((0,3)\)
component, apply the invariant pairing, and integrate the resulting
\((3,3)\)-form.
\end{definition}

\begin{proposition}[Euler--Lagrange equation]
On the ordinary field \(A\in\Omega^{0,1}(X,\g)\),
\[
  I_{\hCS}(A)
  =
  \int_X \Omega_X\wedge
  \left(
    \frac{1}{2}\langle A,\dbar A\rangle_\g
    +
    \frac{1}{6}\langle A,[A,A]\rangle_\g
  \right)
\]
has equation of motion
\[
  F_A^{0,2}:=\dbar A+\frac{1}{2}[A,A]=0.
\]
\end{proposition}

\begin{proof}
For compactly supported variations,
\[
  \delta I_{\hCS}(A)
  =
  \int_X\Omega_X\wedge
  \left\langle
    \delta A,\dbar A+\frac{1}{2}[A,A]
  \right\rangle_\g .
\]
The equality follows from \(\dbar\Omega_X=0\), Stokes' formula with the
support convention above, and invariance of
\(\langle-,-\rangle_\g\). Since \(\delta A\) is arbitrary, the displayed
curvature vanishes.
\end{proof}

\section{Odd symplectic structure and master equation}

\begin{construction}[BV pairing]
Define
\[
  \omega_{\BV}(\alpha,\beta)
  =
  \int_X \Omega_X\wedge \langle \alpha,\beta\rangle_\g .
\]
Only Dolbeault degrees summing to \(3\) contribute. Hence
\[
  \omega_{\BV}(c,c^\vee)
  =
  \int_X \Omega_X\wedge\langle c,c^\vee\rangle_\g,
  \qquad
  \omega_{\BV}(A,A^\vee)
  =
  \int_X \Omega_X\wedge\langle A,A^\vee\rangle_\g .
\]
The ghost numbers of each paired pair sum to \(-1\), so
\(\omega_{\BV}\) has ghost degree \(-1\). On compact \(X\) this is
Serre duality. On \(\C^3\) it is the compact-support/distributional
duality used by the local BV theory.
\end{construction}

\begin{definition}[Hamiltonian BV vector field]
The Hamiltonian vector field of \(I_{\hCS}\) is
\[
  Q_{\BV}\mathcal A
  =
  \dbar\mathcal A+\frac{1}{2}[\mathcal A,\mathcal A].
\]
Equivalently, after projecting to Dolbeault degrees,
\[
\begin{aligned}
  Q_{\BV}c &= \frac{1}{2}[c,c],\\
  Q_{\BV}A &= \dbar c+[A,c],\\
  Q_{\BV}A^\vee &= \dbar A+\frac{1}{2}[A,A]+[A^\vee,c],\\
  Q_{\BV}c^\vee &= \dbar A^\vee+[A,A^\vee]+[c^\vee,c].
\end{aligned}
\]
The signs are the Koszul signs from the graded Dolbeault wedge product.
\end{definition}

\begin{theorem}[Classical master equation]
\[
  \{I_{\hCS},I_{\hCS}\}_{\BV}=0,
  \qquad
  Q_{\BV}^2=0 .
\]
\end{theorem}

\begin{proof}
Let
\[
  F(\mathcal A)=\dbar\mathcal A+\frac{1}{2}[\mathcal A,\mathcal A].
\]
The identity
\[
  \dbar F(\mathcal A)+[\mathcal A,F(\mathcal A)]=0
\]
is the Bianchi identity for the Dolbeault dg Lie algebra
\((\Omega^{0,\bullet}(X,\g),\dbar,[-,-])\). It uses only
\(\dbar^2=0\), the graded Leibniz rule, and Jacobi. Since
\(\delta I_{\hCS}(\mathcal A)=\omega_{\BV}(\delta\mathcal A,F(\mathcal A))\),
the Hamiltonian vector field squares to zero. The same equality is the
classical master equation for the odd Poisson bracket induced by
\(\omega_{\BV}\).
\end{proof}

\section{The flat local propagator on \texorpdfstring{$\C^3$}{C3}}

Fix \(X=\C^3\) with coordinates \(z_1,z_2,z_3\) and
\(\Omega=dz_1\wedge dz_2\wedge dz_3\). Put
\[
  r^2(z,w)=\sum_{i=1}^{3}|z_i-w_i|^2,
  \qquad
  \widehat{d\bar z_k}
  =
  d\bar z_1\wedge\cdots\wedge\widehat{d\bar z_k}
  \wedge\cdots\wedge d\bar z_3 .
\]

\begin{proposition}[Bochner--Martinelli kernel]
On \(\C^3\setminus\Delta\),
\[
  K_{\mathrm{BM}}(z,w)
  =
  \frac{2}{(2\pi i)^3}
  \sum_{k=1}^{3}
  (-1)^{k-1}
  \frac{\overline{z_k-w_k}}{r^6(z,w)}
  \widehat{d\bar z_k}
\]
is a \((0,2)\)-form in \(z\). Away from \(\Delta\) it is
\(\dbar_z\)-closed, and as a current its \(\dbar_z\)-boundary is the
diagonal class. With \(C_\g\in\g\otimes\g\) the inverse tensor to
\(\langle-,-\rangle_\g\),
\[
  P_{\hCS}(z,w)=K_{\mathrm{BM}}(z,w)\otimes C_\g
\]
is the local hCS propagator.
\end{proposition}

\begin{proof}
The coefficient \(2/(2\pi i)^3=(3-1)!/(2\pi i)^3\) is the standard
Bochner--Martinelli normalisation in complex dimension \(3\). Direct
differentiation gives \(\dbar_z K_{\mathrm{BM}}=0\) off the diagonal:
the diagonal terms combine to the radial derivative and the
off-diagonal terms cancel in pairs by antisymmetry of
\(\widehat{d\bar z_k}\). Integration of
\(\Omega_z\wedge K_{\mathrm{BM}}(z,w)\) over a small sphere linking
\(\Delta\) equals \(1\); Stokes' formula gives the diagonal current.
Tensoring with \(C_\g\) inserts the inverse gauge pairing in Wick
contractions.
\end{proof}

\begin{remark}[Free two-point functions]
For abelian \(\g\) the interaction term vanishes and
\(Q_{\BV}=\dbar\). The kernel \(P_{\hCS}\) pairs \(A(z)\) with
\(A^\vee(w)\) and \(c(z)\) with \(c^\vee(w)\), with Dolbeault degrees
complementary to \(3\). For non-abelian \(\g\) the same kernel is the
quadratic part of the perturbative expansion; the cubic vertex is
\(\frac{1}{6}\int\Omega\wedge\langle\mathcal A,[\mathcal A,\mathcal A]\rangle_\g\).
\end{remark}

\section{Quantum master equation}

\begin{definition}[Renormalised QME datum]
A perturbative quantisation of the local theory is a collection of
effective interactions \(I[L]\in\mathcal O_{\mathrm{loc}}(\mathcal E)[[\hbar]]\),
indexed by a length scale \(L>0\), together with a heat-kernel
regularisation of the BV Laplacian \(\Delta_L\), such that
\[
  \left(Q+\hbar\Delta_L\right)\exp(I[L]/\hbar)=0
\]
and the \(I[L]\) satisfy Costello's renormalisation-group flow between
scales.
\end{definition}

\begin{proposition}[Local anomaly class in dimension three]
For six-dimensional hCS the one-loop obstruction to the QME is governed
by the invariant quartic polynomial
\[
  \theta_{\g}\in\left(\Sym^4\g^\vee\right)^\g,
  \qquad
  \theta_{\g}(x)=\operatorname{Tr}_{\g}\bigl(\operatorname{ad}_x^4\bigr)
\]
for the adjoint gauge field contribution, with the evident signed sum
when extra matter fields are present. A local QME solution on the flat
formal chart requires a trivialisation of this class in local Lie
algebra cohomology.
\end{proposition}

\begin{proof}
The only local gauge-invariant polynomial of the required cohomological
degree in complex dimension \(3\) is degree \(4\) in the curvature. The
one-loop wheel graph has four external gauge-field insertions, and its
Lie-algebra factor is the symmetrised trace of four adjoint actions.
Costello's obstruction complex identifies this local functional with
the class \(\theta_\g\in(\Sym^4\g^\vee)^\g\). Vanishing of the class, or
an explicit counterterm or matter contribution with opposite class,
turns the first QME obstruction into a coboundary.
\end{proof}

\section{Compact and global quantisation}

\begin{definition}[Compact hCS quantisation datum]
Let \(X\) be a compact Calabi--Yau threefold. A compact hCS
quantisation datum consists of:
\begin{enumerate}
\item the holomorphic volume form \(\Omega_X\) and an orientation of the
  determinant line of the elliptic BV complex;
\item a Hermitian metric, gauge-fixing operator, heat kernel, and
  propagator compatible with \(\dbar\);
\item renormalised counterterms \(I_k[L]\) satisfying the
  renormalisation-group equation;
\item a trivialisation of the global QME anomaly whose local chart
  representative is \(\theta_\g\);
\item Cech or factorisation descent for observables over a Weiss cover,
  including the boundary or TCFT descent data used by defects;
\item a completed observable complex, for example a nuclear
  Fréchet/continuous-dual completion before formal power series in
  \(\hbar\).
\end{enumerate}
\end{definition}

\begin{proposition}[Scope of the compact claim]
Given the datum above, the renormalised observables
\[
  \Obs^q_{\hCS}(X,\g)
\]
form a global factorisation algebra satisfying the QME. Projecting or
specialising a CY threefold output along the CY-to-chiral functor gives
an \(E_1\)-chiral algebra. \(E_2\)-structure belongs to a centre,
double, module, or enhancement layer.
\end{proposition}

\begin{proof}
The QME supplies a square-zero quantum differential
\(Q+\hbar\Delta_L+\{I[L],-\}_L\) at each scale. The
renormalisation-group equation identifies the complexes at different
scales. The descent datum glues the local factorisation algebras, while
the determinant-line orientation supplies the integration orientation
for BV pushforwards. The final sentence is the \(d\geq3\) operadic
scope convention for \(\Phi_d\).
\end{proof}

\section{Hall and chiral anchors}

\begin{proposition}[The affine toric chart]
For the Jordan quiver with potential attached to the affine toric
Calabi--Yau chart,
\[
  \CoHA(\C^3)=Y^+(\widehat{\mathfrak{gl}}_1).
\]
The full \(\mathcal W_{1+\infty}\) object enters after Drinfeld double,
centre, completion, and representation-theoretic evaluation.
\end{proposition}

\begin{proof}
The critical CoHA of \(\C^3\) is the positive Hall half generated by
the Jordan-quiver creation operators. The double adds the negative
half and Cartan part; completion and evaluation identify the resulting
vacuum representation with the corresponding
\(\mathcal W_{1+\infty}\) module.
\end{proof}

\begin{proposition}[Compact \texorpdfstring{$K3\times E$}{K3 x E} constants]
The compact \(K3\times E\) comparison uses distinct invariants:
\[
  \kcat(K3\times E)=0,\qquad
  \kch(K3\times E)=0,\qquad
  \kch^{\mathrm{Heis}}(K3\times E)=3,
\]
\[
  \kBKM(\Delta_5)=5,\qquad
  \kfib(K3)=24.
\]
\end{proposition}

\begin{proof}
The first equality is Kunneth multiplicativity for
\(\chi(\mathcal O_X)\), using \(\chi(\mathcal O_E)=0\). The compact
Hodge supertrace \(\kch(K3\times E)\) also vanishes in the CY threefold
layer. The value \(3\) belongs to the Heisenberg--Mukai specialisation,
the value \(5\) is the Borcherds weight of Gritsenko's \(\Delta_5\),
and \(24\) is the Mukai-lattice rank of the K3 fibre.
\end{proof}

\section{References used by the construction}

The local BV construction is the holomorphic Chern--Simons example in
Costello--Gwilliam, \emph{Factorization Algebras in Quantum Field
Theory}, Vol. II, Section 5. The heat-kernel and holomorphic
renormalisation framework is Costello--Li, arXiv:1606.00365. The
Bochner--Martinelli kernel is the standard Cauchy--Fantappie kernel in
complex dimension \(3\). The Hall anchor
\(\CoHA(\C^3)=Y^+(\widehat{\mathfrak{gl}}_1)\) is the
Kontsevich--Soibelman and Schiffmann--Vasserot positive-half
identification.

\end{document}
